\section{Cyclic knot-group certificates and adaptive Whitehead descent}
\label{sec:group-certificates}

\subsection{A positive certificate from the full fundamental group}

The optional \code{--group} stage provides a new sufficient test for
unknottedness.  It simplifies a Wirtinger presentation by exact changes of
presentation, records every nontrivial step, and independently replays that
record.  A successful trace ends with precisely one generator and no
nonempty relators.  It therefore proves that the knot group is isomorphic
to $\mathbb Z$.  For classical knots in $S^3$, this implies that the knot is
the unknot, by the classical consequence of the Loop Theorem
\cite{papakyriakopoulos}.  This criterion concerns the whole fundamental
group: its abelianization alone is $\mathbb Z$ for every knot and cannot
supply this conclusion.

Here is the topological implication explicitly.  Write $X$ for the compact
knot exterior.  If $\pi_1(X)=\mathbb Z$, the map from the boundary torus
group $\mathbb Z^2$ cannot be injective.  The Loop Theorem supplies an
embedded compressing disc.  The exterior is irreducible: an embedded
sphere bounds two balls in $S^3$, and the ball on the side away from the
connected knot lies in the exterior.  Compressing its torus boundary along
the disc gives a sphere; irreducibility then identifies the exterior with
a solid torus.  Equivalently, the compression slope is the longitude,
which bounds a disc and gives a spanning disc for the knot after adjoining
the collar annulus.  The input contract is essential: our PD loader
validates a classical, single-component knot diagram.  No conclusion for
virtual diagrams or arbitrary finitely presented groups is asserted.

\subsection{A deterministic presentation and reversible eliminations}

For each PD row $(a,b,c,d)$, join the labels $b,d$ along the overpass.
The connected classes of arc labels are the Wirtinger generators.  Name
these classes by increasing minimum arc label; this makes certificate
indices independent of the producer's union--find root choices.  Orient
the knot by following its component from dart zero.  At a crossing let
$u,v$ denote the incoming and outgoing underpass generators and let $o$
be its overpass generator.  Our relator is
\[
 o u o^{-1}v^{-1}\quad\hbox{at a positive crossing},\qquad
 o v o^{-1}u^{-1}\quad\hbox{at a negative crossing}.
\]
These choices give the same signed Fox rows as the established Alexander
convention.  We retain all crossing relators, including the redundant
Wirtinger relator, and freely and cyclically reduce each one.  Free
reduction changes no word in the free group; cyclic reduction conjugates
the relator and hence preserves its normal closure.

Suppose a relator contains a generator $g$ exactly once, with sign
$\epsilon\in\{1,-1\}$.  Write it as $P g^{\epsilon} S$ where neither
$P$ nor $S$ contains $g^{\pm1}$.  Solve the defining relation:
\[
 g=P^{-1}S^{-1}\quad(\epsilon=1),\qquad
 g=SP\quad(\epsilon=-1).
\]
Substitute this word, and its inverse for $g^{-1}$, into every other
relator.  Delete this one defining relation and the generator.  This is a
Tietze elimination, giving an isomorphic presented group.  In particular,
it is not the operation of imposing an extra relation.  The trace records
only the original relation index and generator; replay derives the
replacement itself.  Empty slots retain their original indices.

The producer prefers an eligible pair minimizing
$(|R|-2)\operatorname{occ}(g)$, followed by relator length, index and generator.
This is a deterministic estimate of expansion cost, not a completeness
argument.  Renumbering or mirroring a diagram can change the greedy path.
Indeed, two of the twelve tested mirror diagrams stall even though every
original succeeds.  This is a concrete reason to retain other decision
procedures in the portfolio.

\subsection{Cut-based changes of free basis}

When no generator can be eliminated, try to shorten the whole relator
list by a Whitehead automorphism.  Fix a signed letter $a$ and a subset
$A$ of the signed alphabet with $a\in A$ and $a^{-1}\notin A$.  Fix
$a^{\pm1}$, and for every other letter $x$ put
\[
 \tau_{A,a}(x)=
 \bigl(a^{-1}\text{ if }x^{-1}\in A\bigr)\;x\;
 \bigl(a\text{ if }x\in A\bigr),
\]
where an absent factor means the empty word.  This preserves inversion and
defines a free-group automorphism.  Its inverse uses multiplier $a^{-1}$
and subset $(A\setminus\{a\})\cup\{a^{-1}\}$.  Applying it to every
relator therefore preserves the isomorphism type of the quotient group.

For a cyclically reduced list of words, add a weighted undirected edge
$\{x,y^{-1}\}$ for each cyclic adjacent pair $xy$, including the last--first
pair of each word.  Add weights across all relators.  For this convention,
the exact change in total cyclic length is
\[
 |\tau_{A,a}(R)|_{\rm cyc}-|R|_{\rm cyc}
   =\operatorname{cut}(A,A^c)-\deg(a).
\]
The cut-capacity approach to Whitehead minimization is developed by
Roig--Ventura--Weil \cite{roig-ventura-weil}; their polynomial algorithm
handles varying free rank, rather than enumerating exponentially many
subsets.  We implement the relevant minimization directly on the sum of
the relators' weighted graphs.  For each signed multiplier, an exact
integer-capacity Edmonds--Karp flow finds a minimum cut separating $a$
from $a^{-1}$.  The best strictly negative length change selects the next
move.  A nonnegative minimum ends this search attempt.

The producer explicitly substitutes the chosen automorphism and checks
that the observed length change equals the cut prediction.  The verifier
needs neither this formula nor the minimum-cut computation: it checks that
the supplied pair is a valid Whitehead automorphism and directly applies
it.  Soundness does not require the proposed move to shorten the words.
Thus an error in search strategy cannot by itself create a false positive
through the certificate interface.

\subsection{A reproducible convention audit of a new preprint}

Kapovich's July 2026 preprint on compressed primitivity
\cite{kapovich-compressed} motivates keeping group words compressed.  It
states polynomial compressed Whitehead-minimality tests at fixed rank,
with further results for compressed primitivity.  It also cautions that
fast compressed operations alone do not bound the number of descent
steps: elementary descent on $a b^{2^m}$ can require exponentially many
steps in the compressed input length.  Fixed-rank hypotheses also need
attention when the rank grows with a knot diagram.

Our audit found a convention mismatch in version~1's equations
(10)--(12), as printed together.  Its automorphism (10) is the one above,
but its graph (11) uses edges $\{x^{-1},y\}$.  Take
\[
 w=abc,\qquad A=\{a,b^{-1}\},\qquad\hbox{multiplier }a.
\]
Literal substitution gives $a a^{-1}bc$, hence cyclic length change $-1$.
The printed graph has edges $\{a^{-1},b\}$, $\{b^{-1},c\}$,
$\{c^{-1},a\}$; the stated cut has weight two and $\deg(a)=1$, predicting
$+1$ via (12).  With our edges $\{x,y^{-1}\}$ the cut has weight zero,
so the predicted change is correctly $-1$.

This local inconsistency is repaired by inverting the graph's vertex
labels, or by transforming the associated characteristic pair
consistently.  It is not a disproof of the paper's compressed-primitivity
claims: minimization over all pairs can survive such a relabelling.  Our
regression test records the three-letter example and checks 15,552
short-word/characteristic-pair combinations against an independent
four-case substitution implementation.  The production code currently
uses explicit word lists.  It does not implement compressed free-group
operations or inherit a compressed-input complexity guarantee.

\subsection{Independent replay and local resource accounting}

The certificate contains its version, method, positive verdict, exact
normalized input PD, move list, and surviving generator.  The checker
reconstructs the presentation by graph traversal rather than union--find,
reconstructs orientation from paired darts rather than cached Diagram
orientation, and performs word cancellation with a separate deque-based
implementation.  It checks unique occurrence before an elimination and
valid alphabet membership and multiplier conditions before a Whitehead
move.  After every substitution it reduces every remaining relator.
Acceptance requires exactly the declared surviving generator and an
entirely empty relator list.  All relators are accounted for; neither the
stored verdict nor a claim of rank one is accepted on its own.

The search and checker share a default wall allowance of 0.05 seconds per
attempt.  Each separately has a two-million-unit cooperative work budget
and a 200,000-letter expansion ceiling.  Work units conservatively charge
letter processing in blocks, graph scans and flow exploration; they are
not a hardware-independent instruction count.  Substitution sizes are
checked before allocating expanded words.  The Whitehead producer uses a
conservative bound of three times the current length.  A trace that is
mathematically valid can therefore be rejected by a resource ceiling.
Only a fully completed and checked trace supplies a verdict.

Local exhaustion or stalled search returns \code{INCONCLUSIVE}; global
cancellation propagates and the surrounding recognizer returns
\code{UNKNOWN}.  The optional stage runs after cheap invariant
obstructions and before the optional normal-surface worker and late
Reidemeister III/Khovanov stages.  A local failure consumes only its
allowance and resumes the portfolio.  Verification is inside the local
allowance, including the check immediately before acceptance.

\subsection{What can and cannot be bounded}

Let $n$ be the crossing number, and suppose every expanded relator list in
one completed search has total length at most $B$.  There are at most
$n-1$ generator eliminations.  Between eliminations each chosen Whitehead
step decreases a nonnegative integer length, so there are at most $B$
such steps in each of at most $n$ phases.  The graph has $O(n)$ vertices
and $O(n^2)$ edges.  Edmonds--Karp takes $O(n^5)$ graph operations per
multiplier and we try $O(n)$ multipliers.  Thus, conservatively, the search
cost is polynomial in $n$ and $B$, bounded by
\[
 O\bigl(n(B+1)(B+n^6)\operatorname{polylog}(nB+2)\bigr)
\]
with straightforward bit accounting for integer capacities.  Replay is
polynomial in explicit trace size and expanded word volume.  This is our
implementation-level conditional bound, not a compressed-word theorem.

Two essential gaps prevent a general fast recognition theorem.  First,
eliminating a generator can substantially expand the other relators; no
polynomial or quasi-polynomial bound on $B$ in terms of $n$ is established.
Second, even unlimited length-decreasing Whitehead search followed by this
greedy elimination rule can stall on an unknot.  A larger Tietze search
would require its own completeness and quantitative bounds.  Compressed
words could address representation growth, but would not alone supply
those missing search bounds.  The bounded optional stage preserves the
existing complete fallback and does not improve the general exponential
upper bound proved for our built-in decision pipeline.

\subsection{Whole-query measurements and regression tests}

\code{fast/benchmark\_group.py} measures 21 inputs in seven randomized
A/A/B rounds after one excluded warmup: 441 timed whole queries.
Every query starts with fresh validated PD data and carries no source-braid
provenance.  The two control arms use the default pipeline.  The third
enables the group stage, including independent replay.  Every arm has a
0.25-second global allowance and a 50,000-object Khovanov ceiling.  Raw
samples, resulting certificates and environment metadata are recorded in
\code{fast/results/group\_20261008.json}.

The group stage certifies all twelve maintained Khovanov-only survivor
inputs in every measured round.  Ten show lower median query costs.  The
two largest gains are $10.90$ to $4.66$ ms and $12.86$ to $3.86$ ms.
Two regress, from $4.62$ to $4.89$ ms and from $3.67$ to $4.69$ ms.
The sum of the twelve median costs is $60.85$ ms for the default control,
$60.56$ ms for the repeated control, and $39.19$ ms with group certificates.
This aggregate describes these twelve selected inputs, not a distribution
of arbitrary knots or an asymptotic law.

An eight-crossing hard unknot and Monster already have cheap default
Reidemeister paths; inserting group search increases their median costs
from $0.66$ to $1.01$ ms and $0.95$ to $1.82$ ms respectively.  Easy
structural and invariant obstructions generally bypass the new stage.
All arms time out on Gordian under this benchmark's short global cap;
these censored observations are not completion times or speedups.
The standalone group attempt also exhausts its local work allowance on
Gordian.  These results support retaining \code{--group} as opt-in.

Tests cover exhaustive short-word cut identities, minimum-cut versus
exhaustive characteristic-pair optimization, the printed-convention
counterexample, all twelve survivor certificates, mirror behaviour,
70 random small knots against an independent full Khovanov cube, known
nontrivial determinant-one knots, malformed and input-mismatched traces,
independent replay with producer helpers disabled, local exhaustion,
global and verification-stage cancellation, and actual CLI routing.

The integrated suite passes all 626 tests with the optional Regina backend
installed (126.707 seconds); see
\code{synthesis/data/group-integrated-tests.txt}.  The rebuilt article
contains 154 pages.
